\diarytitle{0077}{3次正方行列の行列式と逆行列（余因子展開）}

第 1 行に関して余因子展開する。
\[
\det A = 1\cdot\begin{vmatrix}1 & 3\\ 0 & 1\end{vmatrix}
- 2\cdot\begin{vmatrix}0 & 3\\ 2 & 1\end{vmatrix}
+ 0\cdot\begin{vmatrix}0 & 1\\ 2 & 0\end{vmatrix}
= 1\cdot 1 - 2\cdot(-6) + 0 = 1 + 12 = 13
\]
よって
\[
\boxed{\; \det A = 13 \;}
\]
次に各余因子を計算する。
\[
A_{11}=1,\ A_{12}=6,\ A_{13}=-2,\quad
A_{21}=-2,\ A_{22}=1,\ A_{23}=4,\quad
A_{31}=6,\ A_{32}=-3,\ A_{33}=1
\]
（例: $A_{12} = -\begin{vmatrix}0&3\\2&1\end{vmatrix} = -(0-6)=6$、$A_{23} = -\begin{vmatrix}1&2\\2&0\end{vmatrix} = -(0-4)=4$。）
余因子行列を転置して随伴行列を作ると
\[
\operatorname{adj}A =
\begin{pmatrix}
1 & -2 & 6 \\
6 & 1 & -3 \\
-2 & 4 & 1
\end{pmatrix}
\]
したがって
\[
\boxed{\; A^{-1} = \frac{1}{13}
\begin{pmatrix}
1 & -2 & 6 \\
6 & 1 & -3 \\
-2 & 4 & 1
\end{pmatrix} \;}
\]
検算として $AA^{-1}$ を実際に計算する。$A$ の第 1 行 $(1,2,0)$ と $A^{-1}$ の第 1 列 $(1,6,-2)/13$ の内積は $(1+12+0)/13=1$、第 2 列 $(-2,1,4)/13$ との内積は $(-2+2+0)/13=0$、第 3 列 $(6,-3,1)/13$ との内積は $(6-6+0)/13=0$ となる。同様に第 2 行・第 3 行についても計算すると
\[
AA^{-1} =
\begin{pmatrix}
1 & 0 & 0 \\
0 & 1 & 0 \\
0 & 0 & 1
\end{pmatrix} = I
\]
が確かめられる。
